\section{Related Work}\label{sec:related}
% Course rule: about 1/2 AI technology (II-A) + 1/2 domain know-how (II-B).

% M = Methods, AI side (old to new, simple to complex)
\subsection{AI Technology for Price Forecasting and Agents}
\textit{Agricultural price forecasting.} Peng et
al.~\cite{peng2015developing} built an early crop-price forecasting service on
open data from Taiwan's wholesale markets, comparing ARIMA, partial least
squares, and neural networks. Zhang et al.~\cite{zhang2024short} trained an
LSTM on daily prices of six vegetables from seven Beijing wholesale markets and
reported higher accuracy than several machine-learning baselines. Zhao et
al.~\cite{zhao2025vegetable} found that LLM-based forecasters generally
outperform traditional and deep models, but not consistently across vegetables
and time scales, and proposed a mixture of experts. All of these works forecast
wholesale or market averages.

\textit{Tree ensembles and simple baselines.} Random
forests~\cite{breiman2001random} and histogram-based gradient
boosting~\cite{ke2017lightgbm} are strong, data-efficient learners for tabular
features such as recent price changes. However, in the M4 competition, pure
machine-learning methods performed poorly compared with simple statistical
benchmarks~\cite{makridakis2020m4}; a one-layer linear model can beat
Transformer forecasters~\cite{zeng2023transformers}; and removing the LLM from
LLM-based forecasters often does not reduce accuracy~\cite{tan2024language}.
Reliable evaluation must also prevent future information from leaking into
training~\cite{kaufman2012leakage}.

\textit{LLM agents and tool use.} ReAct interleaves reasoning with tool actions
and reduces hallucination in knowledge-intensive tasks~\cite{yao2023react};
Toolformer teaches a model to call APIs such as a
calculator~\cite{schick2023toolformer}. Hallucination nevertheless remains an
open problem~\cite{huang2025survey}. MCP~\cite{anthropic2024mcp} offers an
open, standard way for an agent host to call tool servers, which FMPA uses to
keep all numeric work outside the LLM.

% M = Methods, domain side
\subsection{Domain Knowledge of Fresh Produce Prices}
\textit{Price stickiness.} Retail prices change infrequently: about half of US
consumer prices last less than 4.3 months~\cite{bils2004some}, regular prices
are even more rigid once temporary sales are excluded~\cite{nakamura2008five},
and grocery retailers keep sticky reference prices~\cite{eichenbaum2011reference}.
A no-change forecast is therefore a natural and strong baseline.

\textit{Price transmission.} For fresh vegetables, wholesale prices lead retail
prices and the pass-through is asymmetric~\cite{ward1982asymmetry}; asymmetric
transmission along food supply chains is common~\cite{meyer2004asymmetric}. A
public market average is thus an imperfect proxy for the price in a particular
store.

\textit{Weather shocks and volatility.} Food prices are volatile, and their
volatility must be measured carefully~\cite{gilbert2010food}. In Taiwan,
typhoon warnings raise demand through expectations and stockpiling while
landfall cuts supply, with effects that differ by crop~\cite{yu2022market}; in
the largest cabbage wholesale market, typhoons raise prices, and repeated
storms make the response start earlier and last longer~\cite{su2025food}.

\textit{Public price data.} Taiwan publishes daily wholesale prices and
volumes~\cite{moa2026trans}, and Taichung City publishes daily average prices
of fruits and vegetables in 13 public retail markets, in NT\$ per catty
(600\,g)~\cite{taichung2026retail}. Units differ across sources (per kilogram,
per catty, per bundle), which motivates explicit unit conversion and validation (M1 in
Section~\ref{sec:method}).

% N = New proposed (gap -> our project)
\subsection{Research Gap}
Table~\ref{tab:summary} summarizes representative works.
\textit{Existing research} forecasts wholesale or market averages with
increasingly complex models, studies LLM agents mainly on language tasks, and
explains price formation in economics. \textit{Limitations:} (i) no forecast is
anchored on the price the consumer actually observes; (ii) few studies evaluate
at fixed calendar horizons against no-change and momentum baselines under
leakage control; and (iii) LLM-based systems may let the language model
produce numbers. \textit{Research gap and our project:} FMPA addresses (i)
with store-price anchoring, Eq.~\eqref{eq:modeA}; (ii) with the evaluation
protocol of Section~\ref{sec:experiments}; and (iii) with an MCP tool boundary
that keeps every number outside the LLM (Fig.~\ref{fig:architecture}).

% O = Organize: Article Summary Table at the end of Section II.
% A two-column table* can only float to the top of the next page in IEEEtran.
% Table I -- 8-column article summary required by the course proposal rules.
% "Limitation" is written relative to the FMPA problem setting.
\begin{table*}[t]
\caption{Article summary of representative related work}
\label{tab:summary}
\centering
\scriptsize
\setlength{\tabcolsep}{3pt}
\renewcommand{\arraystretch}{1.12}
\begin{tabularx}{\textwidth}{@{}l c Y Y Y Y Y Y@{}}
\toprule
Ref. & Year & Problem & AI Method & Domain Knowledge & Dataset & Findings & Limitation \\
\midrule
\cite{peng2015developing} & 2015 & Crop price forecasting service & ARIMA, PLS, ANN & Wholesale market open data & Daily prices of four crops, Taiwan & PLS suits short, ANN long horizons & Wholesale only; no store anchor \\
\cite{zhang2024short} & 2024 & Short-term vegetable prices & LSTM vs.\ ML models & Wholesale vegetable markets & 6 vegetables, 7 Beijing markets, 2009--2023 & LSTM most accurate ($R^2=0.958$) & Market averages; compared with learned models \\
\cite{zhao2025vegetable} & 2025 & Vegetable price forecasting & Mixture of LLM and deep experts & Vegetables differ in price dynamics & Vegetable price series & LLMs often best, but inconsistent & No store prices; no agent interface \\
\cite{makridakis2020m4} & 2020 & Which forecasting methods work & 61 statistical, ML, hybrid methods & Naive benchmarks as reference & 100,000 series & Pure ML weak vs.\ simple benchmarks & Not food or retail prices \\
\cite{tan2024language} & 2024 & Are LLMs useful for forecasting & Ablation of LLM components & -- & Standard benchmarks & Removing the LLM does not hurt & Not price data \\
\cite{yao2023react} & 2023 & Ungrounded LLM reasoning & Reasoning + tool actions & -- & QA, fact checking, interactive tasks & Tool calls reduce hallucination & No numeric forecasting \\
\cite{nakamura2008five} & 2008 & How often prices change & -- (descriptive statistics) & Sales vs.\ regular prices & US CPI/PPI microdata & Regular prices are rigid & US data; not produce-specific \\
\cite{ward1982asymmetry} & 1982 & Vegetable price linkages & -- (asymmetry tests) & Wholesale leads retail & US fresh vegetables & Asymmetric pass-through & Old; US markets \\
\cite{meyer2004asymmetric} & 2004 & Asymmetric price transmission & -- (survey of econometric tests) & Farm--wholesale--retail chain & Literature & Asymmetry is common & No forecasting \\
\cite{yu2022market} & 2022 & Typhoon effects on prices & -- (econometric decomposition) & Warning raises demand, landfall cuts supply & Taiwan wholesale vegetables, 1996--2014 & Prices rise; effects vary by crop & Wholesale; pre-2015 data \\
\cite{su2025food} & 2025 & Disaster shocks on food markets & -- (event study, DiD) & Expectations, stockpiling, imports & Daily cabbage, Taiwan's largest wholesale market & Repeated typhoons: earlier, longer response & Cabbage only \\
\midrule
FMPA & 2026 & Store-anchored 7/14-day forecast & LLM + MCP tools; RF/HGB & Stickiness, label maturity, unit normalization & 13 Taichung public markets & Models do not beat naive baselines & Public price $\neq$ store price \\
\bottomrule
\end{tabularx}
\end{table*}

